% ===== PRÉAMBULE CANONIQUE — à coller VERBATIM en tête de CHAQUE .tex =====
\documentclass[11pt,a4paper]{article}
\usepackage[utf8]{inputenc}\usepackage[T1]{fontenc}\usepackage{lmodern}
\usepackage[french]{babel}
\usepackage{amsmath,amssymb,mathtools}\usepackage{icomma}
\usepackage[version=4]{mhchem}          % \ce{...} : équations & formules
\usepackage{siunitx}\sisetup{locale=FR} % \qty{}{}, \si{}, \ang{}
\usepackage{chemfig}                    % \chemfig{...} : structures développées
\usepackage{xcolor}\usepackage{colortbl}
\usepackage{tikz}\usepackage{pgfplots}\pgfplotsset{compat=1.18}
\usepackage[most]{tcolorbox}\usepackage{enumitem}\usepackage{array,booktabs}
\usepackage[a4paper,margin=2.2cm]{geometry}
\usetikzlibrary{arrows.meta,calc,angles,quotes,positioning,patterns,backgrounds,shapes.geometric}
\definecolor{primary}{RGB}{222,49,99}\definecolor{primarydark}{RGB}{150,22,55}
\definecolor{defcol}{RGB}{0,114,178}\definecolor{methcol}{RGB}{52,92,148}
\definecolor{excol}{RGB}{0,135,90}\definecolor{attncol}{RGB}{200,40,55}
\tcbset{boxbase/.style={enhanced,breakable,boxrule=0.9pt,arc=2.5pt,
  left=3mm,right=3mm,top=2.3mm,bottom=2.3mm,fonttitle=\bfseries,coltitle=white}}
% --- boîtes TUTORIEL (noms EXACTS, ne pas renommer) ---
\newtcolorbox{cle}[1][]{boxbase,colback=defcol!6,colframe=defcol,
  title={\textsc{À retenir}\ifx\relax#1\relax\relax\else\textmd{ : \itshape #1}\fi}}
\newtcolorbox{methode}[1][]{boxbase,colback=methcol!7,colframe=methcol,sharp corners,
  title={\textsc{Méthode}\ifx\relax#1\relax\relax\else\textmd{ --- \itshape #1}\fi}}
\newtcolorbox{exemplebox}[1][]{boxbase,colback=excol!5,colframe=excol,
  title={\textsc{Exemple}\ifx\relax#1\relax\relax\else\textmd{ : \itshape #1}\fi}}
\newtcolorbox{piege}[1][]{boxbase,colback=attncol!6,colframe=attncol,
  title={\textsc{Piège}\ifx\relax#1\relax\relax\else\textmd{ : \itshape #1}\fi}}
\newtcolorbox{remarque}[1][]{boxbase,colback=black!4,colframe=black!45,
  title={\textsc{Remarque}\ifx\relax#1\relax\relax\else\textmd{ : \itshape #1}\fi}}
% --- boîtes EXERCICE / CORRIGÉ (noms EXACTS, auto-comptées) ---
\newtcolorbox[auto counter]{exo}[1][]{boxbase,colback=primary!4,colframe=primary,
  title={\textsc{Exercice}~\thetcbcounter\ifx\relax#1\relax\relax\else\textmd{ --- \itshape #1}\fi}}
\newtcolorbox[auto counter]{corrige}[1][]{boxbase,colback=excol!4,colframe=excol,
  title={\textsc{Corrigé}~\thetcbcounter\ifx\relax#1\relax\relax\else\textmd{ --- \itshape #1}\fi}}
\newcommand{\R}{\mathbb{R}}\newcommand{\N}{\mathbb{N}}\newcommand{\Z}{\mathbb{Z}}\newcommand{\ds}{\displaystyle}
\begin{document}
% THEME_TITLE: Masse molaire \& analyse fonctionnelle d'une molécule
\sisetup{per-mode=symbol}

Deux compétences très fréquentes au concours se combinent ici : \textbf{calculer la masse molaire} d'une molécule organique à partir de sa formule, et \textbf{recenser toutes ses fonctions} sur une structure réelle.

\begin{cle}[Masse molaire moléculaire]
\[
\boxed{\,M=\sum_i n_i\,M_i\,}
\]
où $n_i$ est le nombre d'atomes de l'élément $i$ et $M_i$ sa masse molaire atomique (lue dans la classification). Valeurs usuelles :
\[
M(\ce{C})=\qty{12}{\gram\per\mole},\ M(\ce{H})=\qty{1}{\gram\per\mole},\ M(\ce{O})=\qty{16}{\gram\per\mole},\ M(\ce{N})=\qty{14}{\gram\per\mole},\ M(\ce{S})=\qty{32}{\gram\per\mole}.
\]
\end{cle}

\begin{piege}[Masse molaire ($\mathrm{g/mol}$) vs masse moléculaire relative (sans unité)]
La \textbf{masse molaire} $M$ s'exprime en \si{\gram\per\mole}. La \textbf{masse moléculaire relative} $M_r$ est \emph{le même nombre mais sans unité}. Au QCM, une réponse «\,$114\ \mathrm{g/mol}$\,» est \textbf{fausse} si l'on demandait la masse moléculaire \emph{relative} (qui ne porte pas d'unité). Toujours vérifier la grandeur demandée.
\end{piege}

\begin{methode}[De la structure à la masse molaire]
\begin{enumerate}[leftmargin=1.6em,itemsep=0.5pt]
  \item Lire la structure et en déduire la \textbf{formule brute} (compter C, H, O, N, S\dots).
  \item Sommer les contributions : $M=n_{\ce{C}}\cdot 12 + n_{\ce{H}}\cdot 1 + n_{\ce{O}}\cdot 16 + \dots$
  \item Applications : nombre de moles $n=\dfrac{m}{M}$ ; nombre de molécules $N=n\cdot N_{\mathrm{A}}$ avec $N_{\mathrm{A}}=\qty{6.02e23}{\per\mole}$ ; pourcentage massique de l'élément $i$ : $\%_i=\dfrac{n_i M_i}{M}\times 100$.
\end{enumerate}
\end{methode}

\begin{cle}[Analyse fonctionnelle complète d'une molécule (capstone)]
Face à une grosse molécule réelle : (1)~repérer chaque carbonyle $\ce{C=O}$ et son voisinage ; (2)~repérer les hétéroatomes isolés ($-\ce{OH}$, $-\ce{O}-$, $-\ce{N}$, $-\ce{SH}$) ; (3)~repérer les $\ce{C=C}$/$\ce{C#C}$ et les cycles ; (4)~\emph{classer} alcools et amines ; (5)~calculer la formule brute et $M$. On relie ainsi toute la fiche.
\end{cle}

\end{document}
